\chapter{安全・リスク}
\section{背景：危害を生む条件を工程の全体で扱う}
事故を個人の不注意という一語で終えると、同じ事象を起こす設計や運用条件が残る。何が動き、何に接触し、どの防護が機能しなかったのかを先に再現する。その条件を設計、施工、使用、保守、情報の引継ぎまで広げて検討することが、別の業務へ教訓を移す出発点である。

\section{危険源と防護：検知後に止めるだけでよいか}
\label{theory:safety:barriers}
危険源とは危害を生じ得る性質や状況である。重い回転体は大きな運動エネルギーを持ち、人が入る位置で挟まれが起こり得る。リスクを考えるには、その危険へどのように接近し、どの程度の被害が起こり得るかを調べる。

本質的な設計で危害を生むエネルギーや接触を減らす方策、隔離や停止等の防護、残る危険についての情報と使用・保守手順を順に検討する。センサは人を検知した瞬間に機械のエネルギーを消せるわけではない。単純化した回転体で慣性モーメントを\(I\)、角速度を\(\omega\)とすれば運動エネルギーは\(I\omega^2/2\)である。一定の制動トルクの大きさを\(\tau\)とし、制御遅れ等を無視すると、停止までの角変位は
\[
 \Delta\theta=\frac{I\omega^2}{2\tau}
\]
となる。重量・形状や速度を変更したら、検知範囲と停止距離を含めて安全の前提を再評価する必要がある。これは事故機の数値再現ではなく、設計変更が停止性能へ与える影響を示すモデルである。

\section{五つの事例を物理的事象と管理条件へ分ける}
\label{theory:safety:cases}
\begin{description}
 \item[低脂肪乳の食中毒] 停電時の菌増殖と毒素産生が、汚染した脱脂粉乳を通じて最終製品へ及んだ。厚労省の原因説明を用い、菌を殺す工程と既に産生した毒素の問題も分ける\parencite{mhlw_milk2002}。原料の異常から出荷停止・回収までの情報経路は、物理的な原因とともに検討する。
 \item[自動回転ドア] 軽量な設計から機能や外観の変更で重量が増すと、停止できるという前提を再評価する必要がある。大学の事例分析は検知・非常停止への依存と慣性の増大を関連付ける\parencite{muroran_ethics_cases}。顧客の要求に応えたことだけで危険の説明責任が消えるわけではない。
 \item[福知山線] 事故報告の直接原因は速度超過で曲線へ進入したことにあり、ブレーキの遅れと運転士管理方法の関与は、記録に基づく推定や可能性として段階を区別する。曲線の速度超過を防ぐATS機能が使用開始されていれば回避できたと推定されるが、列車に全くATSがなかったわけではない\parencite{jtsb_fukuchiyama2007}。
 \item[戸開走行のエレベーター] 扉が開いたままかごが動くと乗降中の人が挟まれる。2006年の事例はかごの上昇による挟まれであり、ケーブル破断による落下とは機構が異なる。保守の情報と独立した防護の必要性を確認する\parencite{caa_elevator2006}。
 \item[笹子トンネル天井板] 接着系アンカーの脱落等を、材料の劣化という一原因だけへ還元しない。設計の荷重・耐久性、施工と品質記録、点検の方法・頻度、情報の継承を結び付けて扱う\parencite{mlit_sasago2013}。
\end{description}

\section{原因調査と責任の判断を分ける}
\label{theory:safety:responsibility}
再発防止の原因調査は、事象を説明し防護を改善するために行う。刑事責任は、対象被告人の注意義務、予見可能性、結果回避等を個別の事実関係の下で判断する。民事、行政、倫理・組織上の責任も同じ結論になるとは限らない。

歴代3社長を対象とした福知山線の最高裁2017年決定は、個別曲線の危険性の認識等を検討し、無罪の第一審を是認した原判断を相当とした\parencite{court_jr2017}。この結論から、会社の安全上の改善が不要だったとは導けない。事故調査の是正提言と刑事上の注意義務を同一の基準へまとめない。

不起訴処分は、刑事裁判での無罪判決とは異なる。平成30年度問14の公式注記は、笹子の不起訴を判決内容の選択肢へ入れたこと等を理由に不適切出題とする。全員正解という採点措置を、全ての事例の記述が正しいという認定へ置換しない。同問の一部の個別判決資料は未確認であり、問題ブロックにその範囲を残す。

\section{教訓を別の設計へ移す}
事故名と正答番号だけを覚えると、条件が変わった問題へ対応できない。故障機構、危険への接近、防護の作動条件、設計変更、保守で継承すべき情報を取り出す。自分の設備で同じ条件が成立し得るかを点検し、対策の効果と残る危険を記録する。以下では危険源から許容性の判断、対策、記録へ至る共通の枠組みを整理する。
\section{水害・土砂災害：情報を事前の行動計画へ結び付ける}
\label{theory:safety:preparedness}
被害は危険な現象だけで決まらず、曝露する人や設備、避難の可能性にも依存する。
ハザードマップで対象災害、想定の条件、被害範囲、避難先と経路を調べ、
当日の情報をどう行動に変えるか事前に考える。
地図に色がないことを、想定していない現象も含め全て安全だという証明にしない。
避難先と経路は家族の移動能力、地域の危険、所要時間を含めて考える。
備蓄・持出品、連絡方法・集合、保険等は、危険時の移動と被害後の生活の異なる課題を支える。

令和元年度再試験問13の5記述は、当時の一般的な備えとして適切である。
前文のAR5は、評価した排出シナリオのもとで21世紀の全球平均気温上昇を予測する
\cite{env_ar5}。
日本の1901〜2017年の大雨日数増加と降水日数減少も、該当観測資料で扱われる\cite{env_rain2017}。
期間を最新の年へ変えず、全国の集計傾向を全地点・毎年の単調増加へ拡張しない。
温暖化の緩和は将来の危険を減らす課題で、適応・避難の備えを不要にするものではない。

\begin{tcolorbox}[title=出題2019年度と制度改正]
原典は避難準備・高齢者等避難開始、避難勧告、避難指示（緊急）という旧名称を使う。
2021年5月20日から避難勧告は廃止し、従来の勧告の段階から避難指示を行う体系へ一本化した
\cite{env_evac2021}。
旧名称を現在の行動案内としてそのまま使わず、当時の正誤判定を保存する。
\end{tcolorbox}

\begin{tcolorbox}[title=2026年の防災気象情報]
気象庁は2026年5月29日から、情報名称に警戒レベルの数字を付す新たな防災気象情報の運用を開始すると公表した
\cite{env_warning2026}。
市町村等が発令する避難情報と、気象庁等が発表する危険度の情報は主体が異なる。
情報と行動との対応を分かりやすくする改正であり、発令を待つだけで必ず安全になるという保証ではない。
\end{tcolorbox}
注意報・警報等の危険度に応じた発表は、全ての段階が常に順番に来るという約束ではない。
急変で発令や受信が間に合わないこともあるため、危険を認識したときは自ら判断して避難行動を取るという原典の趣旨を理解する。
実際の災害時の個別行動は、現在の地域の情報と自分の条件に基づき決める。
\section{目的と不確かさからリスクを定義する}
\label{theory:safety:definition}
リスクは何が起きるかだけでなく、何を達成したいかに依存する。生産量を増やす設備変更は、
納期を守る機会と、作業者の危険や品質不良の可能性を同時に変え得る。
組織リスクの広い定義は目的への不確かさの影響を扱い、期待からの正負の逸脱を含む。
一方、安全評価では人の危害等の負の側面に焦点を当てる。
この範囲の違いを、互いに矛盾する定義だとしない。
ISOの公開説明も、組織の目的、機会・脅威、資源配分とリスク対応を結び付ける\cite{saf_iso31000}。

\section{危険源・曝露・結果を分ける}
回転する歯車という危険源があっても、閉鎖された領域に人が接近しなければ、
想定する挟まれ事故の経路は成立しない。清掃のためガードを外すと、その経路が開く。
危険源の性質、接触や曝露の機会、防護の状態、起こり得る結果を一つずつ取り出すと、
設計変更、隔離、手順、保護具のどれに作用する対策かが明確になる。
正常運転だけでなく、起動停止、清掃、修理、故障、廃棄でも同じ分解を行う。

\section{確率と損失の積は何を表すか}
\label{theory:safety:quantification}
一つの評価期間に、危害が生じる確率を\(p\)、生じたときの損失を\(L\)とし、
生じないときの損失をゼロとする二状態モデルでは、期待損失は
\[
 \mathbb E[\text{損失}]=pL
\]
である。結果が複数あり、互いに排他的で網羅的な場合には
\[
 \mathbb E[\text{損失}]=\sum_i p_iL_i,\qquad \sum_i p_i=1
\]
と書ける。事象を重複して足さず、期間と曝露集団、損失の単位をそろえる。
これは不確かな結果を平均で評価するモデルであり、最大損失、被害の分布、不公平、
回復不能な危害を全て一つの値だけで表現するものではない。
同じ期待損失でも、稀な大事故と頻繁な軽傷に同じ判断を下すとは限らない。

\section{順序尺度の掛け算と確率を混同しない}
重篤度を1〜5、発生しやすさを1〜5と分類して積を取る手法では、
数値は通常、厳密な損失額や確率ではない。
重篤度4を2のちょうど二倍の被害と解釈したり、分類2を分類1の二倍の確率と
解釈したりするには、別の尺度の根拠が必要である。
行列の評価、加算スコア、段階分岐も可能であり、全てのリスクを確率と損失の積で
定量表現しなければならないという要求にはならない\cite{saf_ra2006}。
手法は対象と利用できる情報に合わせ、重大な危害が平均値に隠れることも点検する。

\section{知識不足とばらつきは異なる}
寸法や故障率をまだ調べていないため分からない状態と、
同じ条件でも故障時点にばらつきがある状態は区別する。
図面、試験、保守記録や外部の事故情報は前者を減らし得る。
ただしデータを増やせば全ての自然変動や未知の事故経路がなくなるわけではない。
既知であることだけでリスクがゼロになるともいえない。
既知の危険な確率事象は依然として安全上のリスクである。

\section{安全と許容性は条件を含む判断である}
\label{theory:safety:acceptability}
安全を許容できないリスクがない状態と捉えると、ゼロリスクであるという証明とは違う。
誰が、どの期間に、どんな危害に曝露されるか、利用の目的と条件、社会の価値、
新たな知見と利用可能な技術を明示する必要がある。
評価基準を一度設定したことは、永久に改訂してはならない理由にならない。
また、限られた設計外力を超える事象があり得ると認識することは、
その被害を社会が無条件に受け入れるべきだという結論ではない。

\section{見積り・分析・評価・対応を分ける}
\label{theory:safety:cycle}
特定は何が起こり得るかを見つける段階である。
見積り・分析は原因、結果、起こりやすさや防護の効き方を調べる段階である。
評価はその情報を基準と比較し、対応の必要性や優先度を判断する段階である。
対応ではリスクを変える方策を選び実施する。
組織リスクの図でアセスメントを特定・分析・評価の三段階として囲んでいても、
職場の実務資料が措置の検討・記録まで含む一連の手法を説明することと直ちに矛盾しない。
設問が採用する用語と囲みの範囲を先に読む。

\section{対策後に戻る理由：残留と新たな危険}
対策の採用後も同じ危険の大きさを再評価する。
遮蔽や囲いを付けたことで温度が上がる、非常停止が通常操作の近くにある、
保守が難しく防護を外したくなる等、新たな経路を生むこともある。
残留リスクとは低減措置後に残るリスクである。
低減したはずという期待だけで完了にせず、許容性と措置の妥当性を確認し、
条件、実施者、判断根拠、残る危険を記録する。
評価書を書くだけでは、接触の可能性や危険源のエネルギーは下がらない。

\section{低減策には優先順位がある}
\label{theory:safety:hierarchy}
製品設計の三段階は、本質的安全設計、防護、使用上の情報という順序で理解する。
職場の2006年指針は、設計・計画での除去低減、工学的対策、管理的対策、
個人用保護具の順を示す\cite{saf_ra2006}。
二つの列挙は用途と粒度が違うため、そのまま同じ四項へ押し込まない。
先に危険な作業をなくす、より危険の小さい材料に替える、動力を減らす等を検討する。
ガード、隔離、排気等で補完し、作業規則と教育、必要な保護具を整える。
人間の注意に依存する方策だけで、上位の検討を省略しない。

\section{ALARPの三つの水準}
\label{theory:safety:alarp}
ALARPは合理的に実行可能な限り低くする原理である。
次の表は、英国議会の報告による許容性の枠組みを独自に整理したものである\cite{saf_parliament2006}。
図形の幅を確率や損失の値と誤読しないよう、数値目盛は付けない。
\begin{center}
\begin{tabular}{p{0.22\linewidth}p{0.65\linewidth}}
\hline
水準 & 判断の意味\\
\hline
許容できない & 低減又は活動の中止を必要とする高い水準である。\\
条件付きで許容 & 合理的に可能な追加低減を行った根拠を要する。\\
広く受容される & 十分低い水準であることと、その状態の維持を検討する。\\
\hline
\end{tabular}
\end{center}
中間へ入っただけで無条件に許容されるわけではない。
下の水準だけを「安全」の唯一の意味とすると、条件付き許容との整合性を失う。
図の説明、規格の適用水準、各国の法的な低減義務を同じ命題へまとめない。

\section{著しい不均衡を予算超過へ置換しない}
追加措置で避ける危害と、その措置に必要な費用、時間、困難等の負担を比較する。
事故の賠償金だけでは回避の便益を表さず、当初事業予算はこの比較の基準とは別である。
例えば一万円でも予算超過なら省略してよいという規則はALARPから導けない。
著しい不均衡を\(C>B\)という単純な大小判定へ置換しない。
ここで\(C\)を費用、\(B\)を同じ尺度で評価できた低減便益としても、
不確かさ、危害の重大性、適用制度の判断が必要である。
あらゆる危険へ使える一定の倍率を、本章で根拠なく作らない。
2006年指針も重大危害と暫定措置の条件を含むため、費用だけで放置を正当化しない。

\section{ALARAは防護の最適化である}
\label{theory:safety:alara}
放射線防護のALARAは合理的に達成可能な限り低くする原理である。
ICRP103第203項の最適化は、線量だけでなく、被ばくの可能性と人数を含み、
経済的・社会的条件を考慮する\cite{saf_icrpaedia}。
防護の改善は目的であるが、線量の数学的な絶対最小値だけを、
生活や社会の他の条件を無視して求める原理ではない。
個人に負担が偏らないような制約も必要である。
食品汚染物質でも国内の分布・摂取等を調べ、国際基準と健康への影響を検討する際に
ALARAが使われてきた\cite{saf_food2023}。
放射線と食品で、同じ定量指標をそのまま流用するわけではない。

\section{低線量モデルと影響の実証を区別する}
ICRPの2007年勧告は確率的影響について低線量の線形応答を防護の基礎に置く
\cite{saf_icrp103}。
モデルの簡略形を\(r(D)=kD\)と書くと、\(D\)は対象とする線量、
\(r\)はその線量に応じた増分リスク、\(k\)は適用する集団・モデルの係数である。
これは生物学的な全影響にしきい値がないという定理ではない。
線量、線量率、影響の種類、集団差と推定の不確かさを切り分ける。
試験の「しきい値がない」を読むときも、確率的影響の防護モデルという範囲を保持する。

\section{リスク対応と負担の移転}
\label{theory:safety:response}
回避は危険な活動を行わない選択、低減は危害の可能性や程度を変える選択、
移転は契約や保険等で負担を他者へ分担させる選択、保有は残る負担を認識して引き受ける選択である。
保険加入で事故の確率や物理的危害そのものが自動的に消えるわけではない。
複数の対応を組み合わせ、契約相手へ危険を知らせず責任を押し付けることと区別する。
対策を実施したという記録だけでなく、望んだ状態への変化を確認する。

\section{リスク認識には同じ式以上のものがある}
\label{theory:safety:perception}
客観的な推定と本人がどれほど危険だと感じるかは別の問題である。
同じ確率でも、本人が選んだ活動か、回避する権限があるか、近い将来か、
誰に利益と損失が分配されるかで受け止め方が違い得る。
同じ「確率と損失の積」という形式を使っても、その二つの主観的入力や価値が違えば
結論は一致しない。
直近の記憶や報道で想起しやすくなること、同じ情報でも利益・損失の提示枠で
判断が変わり得ることも区別する。
これらを全ての人に常に成立する一つの数値法則にせず、設問の一般化を吟味する。

\section{情報交換と市民の参加}
\label{theory:safety:communication}
専門的に細部を説明できない人も、何が不安か、どんな負担を受けるかを表明できる。
情報の不足と価値観の相違を分け、事実、推定、限界、選択肢を伝え、
質問と意見を受け止める。
食品安全委員会の資料は相互作用と参加を重視する\cite{saf_fsc_psych,saf_fsc_glossary}。
専門家の結論を受け入れさせるだけで過程を成功としない。
分からないことは範囲を示し、未知だから既知の被害も説明不要だとしない。

\section{統括管理と全員参加}
\label{theory:safety:governance}
現場の実態は担当者の知識なしには捉えにくい。
その参加を確保しても、必要な人員や資源、方針、措置の実施を統括する責任は残る。
職長の情報収集と、事業の統括管理者の責務を入れ替えない。
新設、変更、非定常作業、新たな事故情報、経年劣化は再調査の契機となる。
調査と記録は説明責任を支えるが、作成だけで法律上の責任が自動免除されるものではない。

\section{参加の効果と残留リスクの規則}
\label{theory:safety:participation}
職場の調査で何が危険かを共有し、合理的に優先度を決める。
残る危険について守る規則と、その必要な理由が分かるようにする。
全員参加の感受性向上という効果と、専門家が分析して危険度を推定する活動は両方あり得る。
公的資料の五つの効果の表現を問う問題では、その列挙との一致を先に確認する
\cite{saf_ra_effects}。
リストにない表現だから一般的にも必ず偽である、という説明を作らない。

\section{経験比は重災害ゼロの保証ではない}
\label{theory:safety:incident}
小さな事故や危険な経験を集めると、重大事故になる前に経路を修正できる。
1:29:300という伝統的な経験比は、その注意を促す説明に用いられる
\cite{saf_welfare}。
集計の期間、業種、危害の定義や報告率が違えば件数は変わる。
29の軽傷要因を除去すれば重大事故が必ずなくなる、各個人が330件で一度重災害を起こす、
という確定的な因果法則へ変えない。
重災害固有の経路も直接調べる。
報告者を責めると情報が隠れるため、早期の報告・保護・改善と共有を結び付ける。

\section{4Sは状態を維持する活動である}
\label{theory:safety:housekeeping}
整理は不要なものを分け除くこと、整頓は必要なものを安全に取り出せる配置へ置くこと、
清掃は汚れを除き異常を見つけること、清潔は良い状態を維持することである
\cite{saf_4s}。
清掃を洗浄へ、清潔をしつけへ置換する列挙は標準的4Sと違う。
しつけを加える列挙は5Sである。
頭文字だけでなく、通路の障害物、漏油、危険物の保管、故障兆候とどう結び付くかを理解する。

\section{警告・使用情報は利用者へ届いて初めて働く}
\label{theory:safety:warnings}
警告は危険源、起こる危害、従わない結果、避ける行動が理解できる必要がある。
表示の可読性、消えにくさ、見つける位置、利用者の言語と能力、使用期間を検討する。
製造国の言語と英語を書いた事実だけで、全ての使用国の利用者が理解できると保証しない。
合理的に予見できる誤使用、保守、運搬、廃棄にも情報を用意する。
ただしあらゆる想像上の誤使用を同じ範囲に含めず、予見可能性の根拠を明示する。
家電の国内指針の適用範囲と、輸出先の法令や表示要求も別である\cite{saf_aeha}。

\section{法令・規格の時点と対象を保存する}
\label{theory:safety:law}
\label{theory:safety:standards}
\begin{tcolorbox}[title=出題当時の制度と法源]
2006年の一般リスク調査の努力義務と、2016年6月1日施行の対象化学物質の調査義務を分ける
\cite{saf_ra2006,saf_chem2016}。
「化学物質」とあるだけで、全製品・全濃度へ無条件の義務を付加しない。
法律本則、政省令、指針、実務解説では規定の場所と性質が違う。
実施時期の詳細が指針にあることを、全部が法律本則の列挙だと言い換えない。
\end{tcolorbox}
SDSは物性・危険有害性・取扱い等の情報を渡す文書で、調査や教育の材料となる。
法的交付・調査義務には指定対象、濃度や用途等の条件がある\cite{saf_chemqa}。
健康診断の実施と結果報告も別であり、一般の定期等の報告の規模条件と
特殊健診の条件を分ける\cite{saf_healthreport}。

\begin{tcolorbox}[title=規格本文の確認範囲]
ISOの公式ページでGuide51の2014版、ISO31000の2018版の情報を確認した。
Guide51の公開全文の取得は失敗しており、本章を全条項の照合済みとはしない
\cite{saf_guide51,saf_iso31000}。
JSAのQ31010:2012の公開表紙はIEC/ISO31010:2009のリスクアセスメント技法を示す
\cite{saf_jis31010}。
平成28年度前文の規格番号の疑義は原典を保存し、正答番号と分離して追加確認へ残す。
平成24年度の1996年規格史と平成25年度の旧家電指針の版も未確認部分を残す。
\end{tcolorbox}

\section{試験での解き方と統合して学ぶ事項}
まず安全の定義、定量モデル、判断基準、措置、記録のどれを問うか決める。
個数問題では各記述の真偽を記録してから数え、候補番号へ変換する。
組合せ問題では表の列順と○×を画像で確かめる。
特に「変わらない」「必ず積」「優先順位なし」「予算に収まらないだけ」という
全称・省略条件を点検する。
公式正答と独立初回が一致しても、他の肢の曖昧さや前文の誤記は解消しない。
頻度は確認済みの概念統計から集計し、異なる回答形式を同一番号の暗記へ縮約しない。
\section{製造物責任を四つの判定に分ける}
\label{theory:pl:framework}
製造物責任法は、製品の欠陥による生命・身体・財産の被害について、製造業者等の民事上の賠償責任を定める。事故が起きたという一事だけで責任が成立するわけではない。まず対象となる物と責任主体を定め、次に欠陥、損害、因果関係を確認し、最後に免責と期間制限を検討する。この順序なら、対象外という判断を安全対策の不要という判断へ飛躍させずに済む。\autocite{pl_statute,pl_qa}

\subsection{過失の証明と欠陥の証明}
\label{theory:pl:proof}
PL法では製造業者の故意又は過失の立証を責任成立の要件としない。ただし、欠陥、損害、両者の因果関係は別の要件である。「過失を証明する必要がない」から「欠陥さえあれば全損害を賠償する」と推論してはならない。H26問7の5とH29問8の2は、救済の変化を短く述べた文として読む一方、独立した法律説明では省略された要件を補って理解する。

\subsection{その製品だけの損害}
\label{theory:pl:selfdamage}
3条ただし書は、損害が当該製造物についてのみ生じた場合をPL法の賠償対象から除く。テレビが発火してテレビだけが壊れた場合と、住宅へ延焼した場合を対比する。前者でも契約等による請求の可能性は別にあり、保証期間が過ぎたという事情だけで全ての救済が消えるとはいえない。自動二輪車の故障も、他の人又は物に被害があるかを先に確かめる。\autocite{pl_statute,pl_qa}

\section{有体物・動産・製造加工の三条件}
\label{theory:pl:scope}
日本法の対象は製造又は加工された動産である。電気、音響、単体ソフトウェア、サービスは有体物という条件を満たさず、土地や家屋自体は動産ではない。ガスやガソリンは有体物であるため、電気と同じ分類に置かない。水道についても、供給サービスと製造加工された物を区別する必要がある。\autocite{pl_definitions2018,pl_qa}

\subsection{単体ソフトウェアと組込み製品}
\label{theory:pl:software}
ソフトウェア自体が対象外でも、その不具合がロボットという製造物の安全性の欠如として現れ、負傷との因果関係が認められるなら、ロボットの製造業者等の責任が問題となる。H24問7のプリインストールされたOSはソフトウェアそのものを、同問のロボットは製品全体を指す。情報の欠陥を有体製品の欠陥へ結び付ける段階が解法の要点である。

\subsection{不動産に取り付けた動産}
\label{theory:pl:attachment}
製造物の引渡し後に建物へ取り付けたという理由だけで、製造時の欠陥の責任が消えるわけではない。エスカレータ、シャッター、サッシ等について、引渡時の動産性と欠陥、損害との因果関係を検討する。建物本体の施工不良とは異なる。逐条解説はエスカレータからの転落事件の裁判例も紹介している。\autocite{pl_definitions2018}

\subsection{冷凍・冷蔵と加工の区別}
\label{theory:pl:processing}
加工は、物の本質を保ちながら新しい属性を付け加える操作として理解する。単なる切断、冷凍、冷蔵、乾燥は基本的にこの加工へ含めない。一方、加熱、調味、塩漬け、燻製等により新しい属性を加えた食品は対象になり得る。冷凍肉とハムを比べる際には、保存状態の変化と新属性の付与を分ける。植物工場のトマトや養殖魚も、生産技術が高度であるという事情だけでは製造加工された物とはいえない。\autocite{pl_definitions2018}

\paragraph{初回解答を訂正する根拠}
H25問7の初回独立解答は、冷凍冷蔵を加工と扱って不適切6件・候補5とした。逐条解説の物理3頁（印刷4頁）を画像で確認し、この分類を訂正した。オを適切とすると、不適切はア・ウ・エ・カ・キの5件となり候補4に対応する。初回記録は保存し、公式解答との一致を訂正理由の代わりにしない。同問イの「品質上の欠陥」は安全性に関係する範囲で読む必要がある。

\section{欠陥は安全性の不足である}
\label{theory:pl:defect}
欠陥の判断には、製品の特性、通常予見される使用形態、引渡時期等を総合して考慮する。単なる性能の不足や外観の不良と、安全性の不足を区別する。設計、製造、指示警告の各段階に安全性の不足が生じ得る。メーカーが意図した操作だけを検討すればよいわけではなく、通常予見できる誤使用も安全設計の検討に含める。\autocite{pl_qa,pl_statute}

\section{製造業者等と免責の条件}
\label{theory:pl:producer}
業として製造・加工・輸入した者に加え、製造業者としての表示、製造業者と誤認させる表示、実質的製造業者と認められる表示をした者が責任主体となり得る。OEMの販売者は「常に除外」でも「常に該当」でもなく、表示と実質を確認する。単なる販売店と輸入者を一括して対象外としてはならない。\autocite{pl_statute,pl_qa}

\subsection{開発危険の抗弁}
\label{theory:pl:development}
引渡時の科学技術の知見では欠陥を認識できなかったことを製造業者等が証明した場合には、4条の免責が問題となる。担当者個人が知らなかったという説明では足りない。「認識できなくても必ず責任を負う」は免責規定を落とした文である。部品又は原材料には、欠陥が専ら完成品側の設計指示によって生じ、かつ部品側に過失がない場合の別の免責もある。

\section{修理された製品と修理行為を分ける}
\label{theory:pl:repair}
性質の回復・維持を行う修理、修繕、整備というサービス自体は、基本的に製造加工ではない。しかし、修理を経た元の製品が製造物でなくなるわけではない。元の引渡時にあった欠陥か、後の修理で生じた不具合かを区別する。H24問7の「使用者が修理した製品」を対象物の定義として読むと対象性は残り、修理行為や修理起因の不具合について読む場合とは数え方が異なる。この問の初回候補3と公式候補4の差は、解釈疑義として保存する。\autocite{pl_qa,pl_definitions2018}

\subsection{再生品と元の製造業者}
\label{theory:pl:recycle}
使用困難になった物の一部から製造加工した再生品は対象となる。再生した者が責任主体となり得るほか、原材料となった元の製品の引渡時の欠陥が再生品の損害と因果関係を持つなら、元の製造業者の責任も検討する。「最後に再生した者が全ての責任を負う」と一律にはいえない。\autocite{pl_definitions2018}

\section{期間制限は出題時点を確認する}
\label{theory:pl:time}
現行法の短期期間は損害及び賠償義務者を知った時から原則3年、生命又は身体侵害では5年である。長期期間は製造業者等の引渡しから10年であり、消費者の購入日から数えるとは限らない。蓄積性物質や潜伏期間のある損害には起算点の特則もある。H30問9の通常のテレビ火災について、欠陥の証明だけで期間経過を無視できるという文は成立しない。\autocite{pl_statute,pl_qa}

\subsection{2020年4月1日の改正と経過措置}
\label{theory:pl:transition}
改正前は生命身体侵害にも短期3年を用いた。2020年4月1日施行の改正ではこれを5年へ延長し、施行時に旧短期期間が完成していない場合の適用を定めた。長期10年についても時効期間であることを明記した。消費者庁の2018年逐条解説は原則として改正後の規定を説明する資料であるため、刊行年をそのまま当時の施行法の年とみなさない。\autocite{pl_transition2020,pl_time2018}

\section{民法・特別法・海外制度との関係}
\label{theory:pl:otherlaw}
PL法に特別の定めがない責任については民法の規定を用いる。対象外となる事例に全く責任がないとは結論できない。原子炉運転等による原子力損害については原子力損害賠償法の規定を確認し、同法4条3項によるPL法不適用を扱う。また、日本法だけを守れば海外の製造物責任制度も満たすという国際共通化の主張は成立しない。逐条解説にも電気の対象性の違いが説明されている。\autocite{pl_nuclear_statute,pl_statute,pl_definitions2018}

\section{賠償責任の回避を超える製品安全}
\label{theory:pl:csr}
R02問12は事故後の賠償だけを問う問題ではない。法令順守後のリスク低減、消費者の期待を踏まえた基準、事故予防、対話、弱者と将来への配慮、緊急時の被害拡大防止、苦情解決の手段を順に検討する。問題に提示された社会的責任の定義に照らすと、7項目はいずれも安全・安心へ向けた行動であり、不適切0件に対応する。ISO26000原規格全文との照合は未完として残す。\autocite{pl_product_csr2016}

\section{維持管理を点検から次の判断へつなぐ}\label{theory:safety:maintenance}
社会資本の機能を持続させるには、状態を点検し、診断し、必要な措置を選び、
その根拠と結果を記録して次の周期へ反映する。
予防的・計画的な修繕は、劣化が大きくなってから一斉更新する負担を避け、
長期の費用の縮減・平準化と利用者に必要な機能の維持へつなげる。
計画を作れば必ず費用が一定になる、という保証ではない。\cite{eng52_infra2014}

平成26年の行動計画は施設分野別の基準・取組を含む。
施設の重要度、状態、材料、環境と分野ごとの最低基準を踏まえて実施するのであり、
全インフラの頻度と内容を一つへ固定するという意味ではない。
分野内での統一した健全度の尺度と、全施設で同じ点検をするという全称を分ける。
構造諸元、点検履歴、修繕記録、損傷位置・変化を蓄積し、
データが不足している状態を把握することも診断の仕事である。

\subsection{新技術と人材・責任を結び付ける}
ロボット、非破壊検査、ICTで効率・観測範囲を改善できるが、
測定の限界・見落とし・適用可能性を評価する能力が必要になる。
施設分野・業務ごとの必要能力を明確にし、担い手を確保・育成する。
資格の保有だけで全分野の経験と判断を代替しない。\cite{eng52_infra2014}

\subsection{行政責任と市民参加は両立する}
管理者の安全確保の責任を、地域住民・民間の協力を一律禁止する根拠にはしない。
国土交通省の平成12年白書にも、行政との合意に基づく道路・河川の清掃等の協働が記録されている。
誰が何を行い、どこから専門点検・管理者の判断が必要になるかを明確にする。
清掃への参加の事実から危険作業や構造診断を無条件に住民へ委任できるとは推論しない。\cite{eng52_adopt}

\section{化学品のGHS：分類と伝達を合わせる}
\label{theory:safety:ghs}
同じ化学品でも国ごとに警告の方法が異なると、運搬・販売・使用の途中で危険を見落とす。
GHSは危険有害性を共通の基準で分類し、その結果をラベルやSDSで伝える枠組みである。
分類の種類と程度を表すクラス・区分、絵表示、注意喚起語、危険有害性情報と注意書きを対応させる。
国内法の表示事項を追加することと国際的な調和は両立する。
当時のGHS案内と平成24年の指針を参照する。\cite{saf53_ghs_old,saf53_ghs_notice}

分類が伝えるのは危険有害性であり、どの作業でも同じ被害確率という意味ではない。
例えば同じ毒性の液体でも、密閉移送と開放した噴霧では曝露の条件が違う。
容器名と供給者が分かることは、追加情報を問い合わせ、事故時に追跡するためにも必要である。
複数のラベル要素の優先規則を、任意に不都合な危険だけ省略できる規則へ変えない。
新しい知見を得たら情報を更新する。当時のJISの細目と現行の表示義務は、版を分けて確認する。

\section{機械安全：本質的な設計と残る危険の情報}
\label{theory:safety:machine}
信頼性は、指定の条件と期間に要求された機能を果たす能力である。
高い信頼性で意図どおり動く装置でも、目的とする動きが人を傷つける設計なら安全とはいえない。
安全性を調べる際は、使用範囲と寿命段階、危険源、危険へ接近する状況を明確にする。
リスク分析の結果を評価し、必要な低減を繰り返す。\cite{saf53_machine2007}

まず危険を生む形状・運動・エネルギーを設計から減らす。
次にガードや保護装置等で接近と危害を防ぐ。それでも残るリスク、正しい使用・点検・訓練等の情報を提供する。
使用上の情報を落として「二つだけの方策」としない。
警告を増やすことだけで本質的な安全設計の検討を代替しない。
機械の製造側と使用側は各段階で役割を持ち、残留リスクがあるから全て使用者の責任という結論にはならない。

システム安全はハードウェアだけでなく、環境・人的条件を含める。
機能安全は必要な安全機能が正しく働くことによる安全で、通常の制御機能全てを意味しない。
安全工学は使用者と作業者双方の危害を扱い、レジリエンスは予防だけでなく
変化への対応や回復できる能力へ目を向ける。原典の対比表現を、全規格の排他的な定義へ広げない。

\section{リコール：危害防止の是正と賠償の成立を分ける}
\label{theory:safety:recall}
安全を損なう製品を把握したら、対象範囲と危害を調べ、回収、修理、交換、使用停止の連絡等を検討する。
劣化や表示の誤りでも是正が必要になり得る。
医薬品に別の品質管理法があることから、回収がないとは言えない。
PMDAは医薬品の回収情報を公表している。\cite{saf53_medicine_recall}

消安法の対象範囲、薬事等の別制度、行政への報告、購入者への連絡、社告を別の問いへ分ける。
対象が全て追跡されて回収できる場合の社告の扱いは、当時の手引きの細条件を要確認とする。
2010手引きの直接全文取得は成功しておらず、平成25年の条件を現行手引きだけから確定しない。\cite{saf53_recall2010}
回収した事実はPL責任の成立を直ちに示さない。欠陥・損害・因果関係等は別に検討する。
一方、賠償が成立しない可能性を、危害防止へ取り組まない根拠にしない。\cite{pl_statute,pl_qa}

\section{重大製品事故：報告と公表の起点を区別する}
\label{theory:safety:productaccident}
重大事故の代表類型は、死亡、治療に要する期間30日以上の重傷病、後遺障害、CO中毒、火災である。
製造・輸入側は事故の発生を知った段階から所定の期限内に報告する。
原因分析が未完という理由で期限を超えてよいとはしない。\cite{saf53_products2012}
事故の日、事業者が知った日、行政が報告を受理した日は一致するとは限らない。
後年取得した事業者用手引きは、知った日を含む10日を示すが、2018施行府令全文の照合と同一視しない。\cite{saf53_products2023}

行政の公表は事業者の報告と別の段階である。
2013年のフローでは、ガス・石油機器や製品起因が疑われる事故の迅速な公表と、
他の段階の報告受理後原則1週間という扱いを区別する。
「原則」を全案件に一律の法定絶対期限へ変えない。\cite{saf53_productflow2013}
通常有すべき安全性と、あらゆる条件での絶対安全は異なる。
合理的に予見できる使用を検討し、残るリスクを説明して減らすことが課題である。

長期使用の点検制度は経年劣化への対応を支える。
平成29年度の9品目は当時の指定として保持し、2021年等の対象見直しを遡及しない。
指定から外れたことも、劣化や保守の問題が消えるという意味ではない。\cite{saf53_oldnine}

\section{食品リスク：評価・管理・対話をつなぐ}
\label{theory:safety:foodroles}
科学的な健康影響を評価する機能と、許容する条件や施策を決め実施する管理の機能を分ける。
食品安全委員会の旧説明は、科学的な評価と管理行政の独立性を示す。\cite{saf53_food_roles2007}
評価結果だけから全ての政策が一意に決まるとは限らず、実行可能性、費用、便益や他の影響も検討する。
管理者の都合で評価結果を改変することも許されない。

リスクコミュニケーションでは、行政が情報を整理し対話の場を設け、
事業者は生産・流通・販売の情報を共有し、消費者は必要な情報を求め意見を述べ、
科学者は根拠と限界を分かりやすく説明する。
全過程における意見と情報の交換であって、結論を一方向に伝えて納得させるだけではない。\cite{saf_fsc_psych,saf_fsc_glossary}
リスク評価と管理の担当を入れ替えた肢は、この機能分担から判断できる。

\section{品質不正：記録の真正性と構造安全を別に検証する}
\label{theory:safety:quality}
測定記録を流用すれば、設計条件どおり施工したかを第三者が検証する根拠を失う。
記録不正だけで全ての構造物が倒壊すると断定できないが、実際の安全性の迅速な確認と
不正行為・管理過程の是正はともに必要である。
2015年の基礎ぐい報告は、責任体制と支持層到達等の技術判断を扱う。\cite{saf53_pile2015}

元請は施工全体を統括し、専門工事の担当は現場で得る情報から技術的に判断する。
設計図があるから現地の地盤確認を省略するとはしない。
機器が不調になったら、作業の停止や再取得、異常の報告と判断の経路を用意する。
ハードウェアの高度化だけでは、記録を偽る組織の圧力や曖昧な責任を取り除けない。
品質不正は他社からも教訓を得られる。トップの優先順位、部門間の役割、検査の独立性まで調べ、
個人の行為と組織条件の両方へ対応する。\cite{saf53_quality2020}

\section{作業の変更と軽微な危害の例外}
\label{theory:safety:workplace}
材料、設備、作業手順や人員の条件を変更したら、以前の評価の前提を確認する。
2006年の指針は、建設物や設備・材料・手順の変更時に調査し、疲労等の付加的な影響も考慮する。
対象の選定では、明らかに軽微な負傷・疾病しか予想されない平坦な歩行等を除外できる。
解説は軽微を医師の治療を要しない程度とし、過去に偶然軽い結果しか出なかったという条件とは区別する。\cite{saf53_workplace2006}

対策の優先順位と実行可能性を検討し、低減を続ける。
ALARPは、追加の低減による便益に対して費用等が著しく不均衡となるかを検討する原則で、
単なる安い対策の選択ではない。法令の最低基準を無視する理由にもならない。\cite{saf_parliament2006}

\section{ヒューマンエラー：結果の手前にある条件}
\label{theory:safety:human}
不適切な操作を認めても、「不注意」で説明を終えない。
知識・経験、注意、伝達、組織、手順の省略、状況への反射、混乱、錯覚、
心身の能力、疲労、単調さなどから、なぜその操作が起きたかを調べる。
教育資料の12分類はこれらの発生要因を含む。\cite{saf53_human12}
疲労や年齢による機能変化は、誤操作という行為自体と同じ概念ではない。
令和2年の問いは発生要因の分類に対応するため0項目を選ぶが、前文の9割以上という数値は
母集団・集計法が未確認である。全分野の普遍的な事故率として使わない。

対策は作業者への注意だけに限らない。認知しやすい表示、誤操作を防ぐ形状、
異常時の停止、負荷と休息、明確な連絡経路、責任分担と記録を組み合わせる。
原因の分類名を覚えるだけで、具体的な工程や条件を調べる課題を代替しない。

\section{2021行動計画：章と七つの取組を読む}
\label{theory:safety:maintenance2021}
第2次インフラ長寿命化計画は2021年6月18日に公表され、令和3〜7年度を計画期間とする。
必要施策の方向性にある七項目は、個別施設計画、点検診断・修繕更新、予算管理、体制、
新技術、情報基盤、基準類である。現状課題の別章と見出し数を混同しない。\cite{saf53_infra2021}
技術継承が独立した七項目の見出しにないことは、人材の確保や能力の継承が不要という意味ではない。
計画の階層を読み、見出しへの所属と実際の施策内容を区別する。

\section{事業継続：原因より重要事業への影響を調べる}
\label{theory:safety:bcp}
BCPは事業継続計画、BCMはその計画を支える経営・管理の活動である。重要な事業を継続し、又は目標とする期間・水準で復旧させるために、事業への影響、必要な資源、依存関係、代替手段を整理する。自然災害に限らず、感染症、設備事故、サイバー攻撃、取引先の中断等も事業へ影響する\parencite{eth54_bcp2013,eth54_bcp2021}。

復旧が容易な仕事を先に再開すればよいとは限らない。供給停止が顧客、収入、雇用、法令・契約、地域機能へ与える影響と許容時間を評価し、重要業務を選ぶ。バックアップは本体と同じ原因で失われないよう、場所、電源、通信、供給者等の共通の依存を点検する。

\paragraph{計画から実行・改善へ}
経営者が責任を持ち、訓練、点検、資源の確保、事業環境の変化への見直しを進める。計画書を作ったことだけでは対応能力を保証しない。重要業務や依存関係の可視化は平時の改善にもつながる。地域と取引先との協力を含め、組織単独で復旧できると仮定しない。

非常時だから法令や契約がすべて無効になるわけではない。緊急措置には固有の条件があり、権限や責任を確認する。「全て遵守しなくてよい」という包括的免除とは区別する。

\paragraph{普及率は調査年と分母を読む}
平成29年度調査では策定済みの割合が大企業64.0\%、中堅企業31.8\%とされる\parencite{eth54_bcp_survey2018}。これはすべての会社を同じ分母で測った数字ではない。策定中も含む率、企業規模、回答数、調査時点を分ける。令和元年度問題を後に公表された調査で説明しない。

\section{予防的取組：不確実性の中で行動を選ぶ}
\label{theory:safety:precaution}
知識が不完全なことと危険がないことは別である。深刻又は不可逆的な影響が懸念されるときは、科学的確実性がないことだけを理由に合理的な低減策を延期しない。利益、費用、代替案、行動しない場合の損失を比較し、対策の程度と見直しの条件を定める\parencite{eth54_rio,eth54_eea}。

平成25年度問題の十二項目は、無知と不確実性、監視、科学の盲点、学際的学習、実世界条件、利益とリスク、代替策、地域の知識、異なる価値、規制の独立性、制度の障壁、分析だけによる停滞を点検する。これらを考慮することは、科学を捨てることやあらゆる技術を禁止することではない。原典前文の宣言の「批准」という表現も、条約の法的手続と同一にしない。

\section{対話の視点と検証可能性}
\label{theory:safety:riskdialogue}
管理する側の集団的な確率評価と、危険を受ける当事者の判断は異なる。自分で選択できるか、利益と危険の配分は公平か、組織を信頼できるか等もリスク認知に関わる。ハザードとアウトレージを合わせる表現は認知モデルであり、単位の異なる物理量を足す計算式ではない\parencite{eth54_riskcomm}。

専門家の知識の伝達だけで対話が完成するわけではない。中立的な媒介者の支援、権限と情報の非対称性への配慮、異なる根拠の提示、検証可能性が必要になる。分かりやすい説明を作るときにも、情報源や不確かさを隠してはならない。

\subsection{平成23年度 問5}
\label{q:h23-reg-apt-05}

\noindent\textbf{原典：} sources/H23/aptitude.pdf、PDF物理6頁。

\textbf{公式正答：} 4（single）。

\par\textit{原典本文の校正と独自解説は未作成。分類のみの登録項目。}

\subsection{平成23年度 問6}
\label{q:h23-reg-apt-06}

\noindent\textbf{原典：} sources/H23/aptitude.pdf、PDF物理7頁。

\textbf{公式正答：} 4（single）。

\par\textit{原典本文の校正と独自解説は未作成。分類のみの登録項目。}

\subsection{平成23年度 問8}
\label{q:h23-reg-apt-08}

\noindent\textbf{原典：} sources/H23/aptitude.pdf、PDF物理9頁。

\textbf{公式正答：} 5（single）。

\par\textit{原典本文の校正と独自解説は未作成。分類のみの登録項目。}

\subsection{平成24年度 問6}
\label{q:h24-reg-apt-06}

\noindent\textbf{原典：} sources/H24/aptitude.pdf、PDF物理5頁。

\textbf{公式正答：} 3（single）。

\par\textit{原典本文の校正と独自解説は未作成。分類のみの登録項目。}

\subsection{平成24年度 問7}
\label{q:h24-reg-apt-07}

\noindent\textbf{原典：} sources/H24/aptitude.pdf、PDF物理6頁。

\textbf{公式正答：} 4（single）。

\par\textit{原典本文の校正と独自解説は未作成。分類のみの登録項目。}

\subsection{平成24年度 問10}
\label{q:h24-reg-apt-10}

\noindent\textbf{原典：} sources/H24/aptitude.pdf、PDF物理9頁。

\textbf{公式正答：} 1（single）。

\par\textit{原典本文の校正と独自解説は未作成。分類のみの登録項目。}

\subsection{平成24年度 問11}
\label{q:h24-reg-apt-11}

\noindent\textbf{原典：} sources/H24/aptitude.pdf、PDF物理10頁。

\textbf{公式正答：} 3（single）。

\par\textit{原典本文の校正と独自解説は未作成。分類のみの登録項目。}

\subsection{平成25年度 問4}
\label{q:h25-reg-apt-04}

\noindent\textbf{原典：} sources/H25/aptitude.pdf、PDF物理5頁。

\textbf{公式正答：} 5（single）。

\par\textit{原典本文の校正と独自解説は未作成。分類のみの登録項目。}

\subsection{平成25年度 問5}
\label{q:h25-reg-apt-05}

\noindent\textbf{原典：} sources/H25/aptitude.pdf、PDF物理6頁。

\textbf{公式正答：} 3（single）。

\par\textit{原典本文の校正と独自解説は未作成。分類のみの登録項目。}

\subsection{平成25年度 問7}
\label{q:h25-reg-apt-07}

\noindent\textbf{原典：} sources/H25/aptitude.pdf、PDF物理8頁。

\textbf{公式正答：} 4（single）。

\par\textit{原典本文の校正と独自解説は未作成。分類のみの登録項目。}

\subsection{平成25年度 問8}
\label{q:h25-reg-apt-08}

\noindent\textbf{原典：} sources/H25/aptitude.pdf、PDF物理9頁。

\textbf{公式正答：} 2（single）。

\par\textit{原典本文の校正と独自解説は未作成。分類のみの登録項目。}

\subsection{平成25年度 問9}
\label{q:h25-reg-apt-09}

\noindent\textbf{原典：} sources/H25/aptitude.pdf、PDF物理10頁。

\textbf{公式正答：} 5（single）。

\par\textit{原典本文の校正と独自解説は未作成。分類のみの登録項目。}

\subsection{平成26年度 問7}
\label{q:h26-reg-apt-07}

\noindent\textbf{原典：} sources/H26/aptitude.pdf、PDF物理7頁。

\textbf{公式正答：} 4（single）。

\par\textit{原典本文の校正と独自解説は未作成。分類のみの登録項目。}

\subsection{平成26年度 問8}
\label{q:h26-reg-apt-08}

\noindent\textbf{原典：} sources/H26/aptitude.pdf、PDF物理8頁。

\textbf{公式正答：} 3（single）。

\par\textit{原典本文の校正と独自解説は未作成。分類のみの登録項目。}

\subsection{平成26年度 問9}
\label{q:h26-reg-apt-09}

\noindent\textbf{原典：} sources/H26/aptitude.pdf、PDF物理9頁。

\textbf{公式正答：} 4（single）。

\par\textit{原典本文の校正と独自解説は未作成。分類のみの登録項目。}

\subsection{平成26年度 問11}
\label{q:h26-reg-apt-11}

\noindent\textbf{原典：} sources/H26/aptitude.pdf、PDF物理11頁。

\textbf{公式正答：} 1（single）。

\par\textit{原典本文の校正と独自解説は未作成。分類のみの登録項目。}

\subsection{平成26年度 問13}
\label{q:h26-reg-apt-13}

\noindent\textbf{原典：} sources/H26/aptitude.pdf、PDF物理13頁。

\textbf{公式正答：} 5（single）。

\par\textit{原典本文の校正と独自解説は未作成。分類のみの登録項目。}

\subsection{平成26年度 問15}
\label{q:h26-reg-apt-15}

\noindent\textbf{原典：} sources/H26/aptitude.pdf、PDF物理15頁。

\textbf{公式正答：} 3（single）。

\par\textit{原典本文の校正と独自解説は未作成。分類のみの登録項目。}

\subsection{平成27年度 問4}
\label{q:h27-reg-apt-04}

\noindent\textbf{原典：} sources/H27/aptitude.pdf、PDF物理4頁。

\textbf{公式正答：} 3（single）。

\par\textit{原典本文の校正と独自解説は未作成。分類のみの登録項目。}

\subsection{平成27年度 問5}
\label{q:h27-reg-apt-05}

\noindent\textbf{原典：} sources/H27/aptitude.pdf、PDF物理5頁。

\textbf{公式正答：} 5（single）。

\par\textit{原典本文の校正と独自解説は未作成。分類のみの登録項目。}

\subsection{平成27年度 問9}
\label{q:h27-reg-apt-09}

\noindent\textbf{原典：} sources/H27/aptitude.pdf、PDF物理9頁。

\textbf{公式正答：} 3（single）。

\par\textit{原典本文の校正と独自解説は未作成。分類のみの登録項目。}

\subsection{平成27年度 問13}
\label{q:h27-reg-apt-13}

\noindent\textbf{原典：} sources/H27/aptitude.pdf、PDF物理13頁。

\textbf{公式正答：} 2（single）。

\par\textit{原典本文の校正と独自解説は未作成。分類のみの登録項目。}

\subsection{平成27年度 問14}
\label{q:h27-reg-apt-14}

\noindent\textbf{原典：} sources/H27/aptitude.pdf、PDF物理14頁。

\textbf{公式正答：} 5（single）。

\par\textit{原典本文の校正と独自解説は未作成。分類のみの登録項目。}

\subsection{平成28年度 問3}
\label{q:h28-reg-apt-03}

\noindent\textbf{原典：} sources/H28/aptitude.pdf、PDF物理3頁。

\textbf{公式正答：} 1（single）。

\par\textit{原典本文の校正と独自解説は未作成。分類のみの登録項目。}

\subsection{平成28年度 問5}
\label{q:h28-reg-apt-05}

\noindent\textbf{原典：} sources/H28/aptitude.pdf、PDF物理5頁。

\textbf{公式正答：} 4（single）。

\par\textit{原典本文の校正と独自解説は未作成。分類のみの登録項目。}

\subsection{平成28年度 問7}
\label{q:h28-reg-apt-07}

\noindent\textbf{原典：} sources/H28/aptitude.pdf、PDF物理7頁。

\textbf{公式正答：} 5（single）。

\par\textit{原典本文の校正と独自解説は未作成。分類のみの登録項目。}

\subsection{平成28年度 問9}
\label{q:h28-reg-apt-09}

\noindent\textbf{原典：} sources/H28/aptitude.pdf、PDF物理9頁。

\textbf{公式正答：} 5（single）。

\par\textit{原典本文の校正と独自解説は未作成。分類のみの登録項目。}

\subsection{平成28年度 問15}
\label{q:h28-reg-apt-15}

\noindent\textbf{原典：} sources/H28/aptitude.pdf、PDF物理15頁。

\textbf{公式正答：} 3（single）。

\par\textit{原典本文の校正と独自解説は未作成。分類のみの登録項目。}

\subsection{平成29年度 問5}
\label{q:h29-reg-apt-05}

\noindent\textbf{原典：} sources/H29/aptitude.pdf、PDF物理5頁。

\textbf{公式正答：} 3（single）。

\par\textit{原典本文の校正と独自解説は未作成。分類のみの登録項目。}

\subsection{平成29年度 問8}
\label{q:h29-reg-apt-08}

\noindent\textbf{原典：} sources/H29/aptitude.pdf、PDF物理8頁。

\textbf{公式正答：} 5（single）。

\par\textit{原典本文の校正と独自解説は未作成。分類のみの登録項目。}

\subsection{平成29年度 問9}
\label{q:h29-reg-apt-09}

\noindent\textbf{原典：} sources/H29/aptitude.pdf、PDF物理9頁。

\textbf{公式正答：} 3（single）。

\par\textit{原典本文の校正と独自解説は未作成。分類のみの登録項目。}

\subsection{平成29年度 問12}
\label{q:h29-reg-apt-12}

\noindent\textbf{原典：} sources/H29/aptitude.pdf、PDF物理11頁。

\textbf{公式正答：} 3（single）。

\par\textit{原典本文の校正と独自解説は未作成。分類のみの登録項目。}

\subsection{平成30年度 問9}
\label{q:h30-reg-apt-09}

\noindent\textbf{原典：} sources/H30/aptitude.pdf、PDF物理9頁。

\textbf{公式正答：} 3（single）。

\par\textit{原典本文の校正と独自解説は未作成。分類のみの登録項目。}

\subsection{平成30年度 問10}
\label{q:h30-reg-apt-10}

\noindent\textbf{原典：} sources/H30/aptitude.pdf、PDF物理10頁。

\textbf{公式正答：} 5（single）。

\par\textit{原典本文の校正と独自解説は未作成。分類のみの登録項目。}

\subsection{平成30年度 問11}
\label{q:h30-reg-apt-11}

\noindent\textbf{原典：} sources/H30/aptitude.pdf、PDF物理11頁。

\textbf{公式正答：} 3（single）。

\par\textit{原典本文の校正と独自解説は未作成。分類のみの登録項目。}

\subsection{令和元年度 問3}
\label{q:r01-reg-apt-03}

\noindent\textbf{原典：} sources/R01/aptitude.pdf、PDF物理4頁。

\textbf{公式正答：} 4（single）。

\par\textit{原典本文の校正と独自解説は未作成。分類のみの登録項目。}

\subsection{令和元年度 問7}
\label{q:r01-reg-apt-07}

\noindent\textbf{原典：} sources/R01/aptitude.pdf、PDF物理8頁。

\textbf{公式正答：} 5（single）。

\par\textit{原典本文の校正と独自解説は未作成。分類のみの登録項目。}

\subsection{令和元年度 問8}
\label{q:r01-reg-apt-08}

\noindent\textbf{原典：} sources/R01/aptitude.pdf、PDF物理9頁。

\textbf{公式正答：} 2（single）。

\par\textit{原典本文の校正と独自解説は未作成。分類のみの登録項目。}

\subsection{令和元年度 問11}
\label{q:r01-reg-apt-11}

\noindent\textbf{原典：} sources/R01/aptitude.pdf、PDF物理12頁。

\textbf{公式正答：} 2（single）。

\par\textit{原典本文の校正と独自解説は未作成。分類のみの登録項目。}

\subsection{令和元年度 問12}
\label{q:r01-reg-apt-12}

\noindent\textbf{原典：} sources/R01/aptitude.pdf、PDF物理13頁。

\textbf{公式正答：} 1（single）。

\par\textit{原典本文の校正と独自解説は未作成。分類のみの登録項目。}

\subsection{令和元年度 問13}
\label{q:r01-reg-apt-13}

\noindent\textbf{原典：} sources/R01/aptitude.pdf、PDF物理13頁。

\textbf{公式正答：} 4（single）。

\par\textit{原典本文の校正と独自解説は未作成。分類のみの登録項目。}

\subsection{令和元年度再試験 問5}
\label{q:r01-ret-apt-05}

\noindent\textbf{原典：} sources/R01\_retest/aptitude.pdf、PDF物理5頁。

\textbf{公式正答：} 3（single）。

\par\textit{原典本文の校正と独自解説は未作成。分類のみの登録項目。}

\subsection{令和元年度再試験 問7}
\label{q:r01-ret-apt-07}

\noindent\textbf{原典：} sources/R01\_retest/aptitude.pdf、PDF物理7頁。

\textbf{公式正答：} 3（single）。

\par\textit{原典本文の校正と独自解説は未作成。分類のみの登録項目。}

\subsection{令和元年度再試験 問12}
\label{q:r01-ret-apt-12}

\noindent\textbf{原典：} sources/R01\_retest/aptitude.pdf、PDF物理13頁。

\textbf{公式正答：} 2（single）。

\par\textit{原典本文の校正と独自解説は未作成。分類のみの登録項目。}

\subsection{令和2年度 問6}
\label{q:r02-reg-apt-06}

\noindent\textbf{原典：} sources/R02/aptitude.pdf、PDF物理6頁。

\textbf{公式正答：} 5（single）。

\par\textit{原典本文の校正と独自解説は未作成。分類のみの登録項目。}

\subsection{令和2年度 問7}
\label{q:r02-reg-apt-07}

\noindent\textbf{原典：} sources/R02/aptitude.pdf、PDF物理7頁。

\textbf{公式正答：} 4（single）。

\par\textit{原典本文の校正と独自解説は未作成。分類のみの登録項目。}

\subsection{令和2年度 問9}
\label{q:r02-reg-apt-09}

\noindent\textbf{原典：} sources/R02/aptitude.pdf、PDF物理9頁。

\textbf{公式正答：} 3（single）。

\par\textit{原典本文の校正と独自解説は未作成。分類のみの登録項目。}

\subsection{令和2年度 問13}
\label{q:r02-reg-apt-13}

\noindent\textbf{原典：} sources/R02/aptitude.pdf、PDF物理13頁。

\textbf{公式正答：} 4（single）。

\par\textit{原典本文の校正と独自解説は未作成。分類のみの登録項目。}

\subsection{令和3年度 問2}
\label{q:r03-reg-apt-02}

\noindent\textbf{原典：} sources/R03/aptitude.pdf、PDF物理2頁。

\textbf{公式正答：} 4（single）。

\par\textit{原典本文の校正と独自解説は未作成。分類のみの登録項目。}

\subsection{令和3年度 問8}
\label{q:r03-reg-apt-08}

\noindent\textbf{原典：} sources/R03/aptitude.pdf、PDF物理8頁。

\textbf{公式正答：} 3（single）。

\par\textit{原典本文の校正と独自解説は未作成。分類のみの登録項目。}

\subsection{令和3年度 問10}
\label{q:r03-reg-apt-10}

\noindent\textbf{原典：} sources/R03/aptitude.pdf、PDF物理10頁。

\textbf{公式正答：} 3（single）。

\par\textit{原典本文の校正と独自解説は未作成。分類のみの登録項目。}

\subsection{令和3年度 問15}
\label{q:r03-reg-apt-15}

\noindent\textbf{原典：} sources/R03/aptitude.pdf、PDF物理15頁。

\textbf{公式正答：} 3（single）。

\par\textit{原典本文の校正と独自解説は未作成。分類のみの登録項目。}

\subsection{令和4年度 問5}
\label{q:r04-reg-apt-05}

\noindent\textbf{原典：} sources/R04/aptitude.pdf、PDF物理4頁。

\textbf{公式正答：} 4（single）。

\par\textit{原典本文の校正と独自解説は未作成。分類のみの登録項目。}

\subsection{令和4年度 問6}
\label{q:r04-reg-apt-06}

\noindent\textbf{原典：} sources/R04/aptitude.pdf、PDF物理5頁。

\textbf{公式正答：} 5（single）。

\par\textit{原典本文の校正と独自解説は未作成。分類のみの登録項目。}

\subsection{令和4年度 問7}
\label{q:r04-reg-apt-07}

\noindent\textbf{原典：} sources/R04/aptitude.pdf、PDF物理6頁。

\textbf{公式正答：} 3（single）。

\par\textit{原典本文の校正と独自解説は未作成。分類のみの登録項目。}

\subsection{令和4年度 問8}
\label{q:r04-reg-apt-08}

\noindent\textbf{原典：} sources/R04/aptitude.pdf、PDF物理7頁。

\textbf{公式正答：} 2（single）。

\par\textit{原典本文の校正と独自解説は未作成。分類のみの登録項目。}

\subsection{令和4年度 問11}
\label{q:r04-reg-apt-11}

\noindent\textbf{原典：} sources/R04/aptitude.pdf、PDF物理10頁。

\textbf{公式正答：} 4（single）。

\par\textit{原典本文の校正と独自解説は未作成。分類のみの登録項目。}

\subsection{令和5年度 問3}
\label{q:r05-reg-apt-03}

\noindent\textbf{原典：} sources/R05/aptitude.pdf、PDF物理3頁。

\textbf{公式正答：} 2（single）。

\par\textit{原典本文の校正と独自解説は未作成。分類のみの登録項目。}

\subsection{令和5年度 問6}
\label{q:r05-reg-apt-06}

\noindent\textbf{原典：} sources/R05/aptitude.pdf、PDF物理6頁。

\textbf{公式正答：} 3（single）。

\par\textit{原典本文の校正と独自解説は未作成。分類のみの登録項目。}

\subsection{令和5年度 問8}
\label{q:r05-reg-apt-08}

\noindent\textbf{原典：} sources/R05/aptitude.pdf、PDF物理8頁。

\textbf{公式正答：} 2（single）。

\par\textit{原典本文の校正と独自解説は未作成。分類のみの登録項目。}

\subsection{令和5年度 問12}
\label{q:r05-reg-apt-12}

\noindent\textbf{原典：} sources/R05/aptitude.pdf、PDF物理12頁。

\textbf{公式正答：} 4（single）。

\par\textit{原典本文の校正と独自解説は未作成。分類のみの登録項目。}

\subsection{令和5年度 問14}
\label{q:r05-reg-apt-14}

\noindent\textbf{原典：} sources/R05/aptitude.pdf、PDF物理14頁。

\textbf{公式正答：} 3（single）。

\par\textit{原典本文の校正と独自解説は未作成。分類のみの登録項目。}

\subsection{令和6年度 問7}
\label{q:r06-reg-apt-07}

\noindent\textbf{原典：} sources/R06/aptitude.pdf、PDF物理7頁。

\textbf{公式正答：} 4（single）。

\par\textit{原典本文の校正と独自解説は未作成。分類のみの登録項目。}

\subsection{令和6年度 問13}
\label{q:r06-reg-apt-13}

\noindent\textbf{原典：} sources/R06/aptitude.pdf、PDF物理13頁。

\textbf{公式正答：} 5（single）。

\par\textit{原典本文の校正と独自解説は未作成。分類のみの登録項目。}

\subsection{令和6年度 問14}
\label{q:r06-reg-apt-14}

\noindent\textbf{原典：} sources/R06/aptitude.pdf、PDF物理14頁。

\textbf{公式正答：} 3（single）。

\par\textit{原典本文の校正と独自解説は未作成。分類のみの登録項目。}

\subsection{令和7年度 問12}
\label{q:r07-reg-apt-12}

\noindent\textbf{原典：} sources/R07/aptitude.pdf、PDF物理12頁。

\textbf{公式正答：} 4（single）。

\par\textit{原典本文の校正と独自解説は未作成。分類のみの登録項目。}

\subsection{令和7年度 問13}
\label{q:r07-reg-apt-13}

\noindent\textbf{原典：} sources/R07/aptitude.pdf、PDF物理13頁。

\textbf{公式正答：} 5（single）。

\par\textit{原典本文の校正と独自解説は未作成。分類のみの登録項目。}

\subsection{令和7年度 問14}
\label{q:r07-reg-apt-14}

\noindent\textbf{原典：} sources/R07/aptitude.pdf、PDF物理14頁。

\textbf{公式正答：} 2（single）。

\par\textit{原典本文の校正と独自解説は未作成。分類のみの登録項目。}

\subsection{令和7年度 問15}
\label{q:r07-reg-apt-15}

\noindent\textbf{原典：} sources/R07/aptitude.pdf、PDF物理15頁。

\textbf{公式正答：} 2（single）。

\par\textit{原典本文の校正と独自解説は未作成。分類のみの登録項目。}

